\documentclass[11pt,a4paper]{article}
\usepackage[margin=25mm]{geometry}
\usepackage{amsmath,amssymb,amsthm}
\setlength{\emergencystretch}{2em}
\newtheorem{lemma}{Lemma}
\newtheorem{theorem}[lemma]{Theorem}
\newcommand{\gen}[1]{\langle#1\rangle}
\newcommand{\Alt}{\operatorname{Alt}}
\newcommand{\Sym}{\operatorname{Sym}}
\title{The endpoint cut with first shared label three}
\author{Human Compression project: L044 continuation}
\date{}
\begin{document}\maketitle

\begin{lemma}[An anchored path with at most one extra point]\label{path}
Let $T=\gen{t_0,\ldots,t_{h-1}}$ be transitive on $V$, where
$h\geq3$, the $t_i$ are involutions, and distant indices commute.
Suppose $a\in V$ is fixed by $t_1,\ldots,t_{h-1}$ and
$t_{h-1}\ne1$. If $|V|\leq h+2$, there are distinct points
$a_0=a,a_1,\ldots,a_h$ such that $t_i$ exchanges $a_i,a_{i+1}$
and fixes $a_0,\ldots,a_{i-1}$.
If $|V|=h+1$, then
\[
 T=\Sym(V),\qquad
 \gen{t_1,\ldots,t_{h-1}}=\Sym(V\setminus\{a\}).        \tag{1}
\]
If $|V|=h+2$, the $t_i$ cannot all be odd permutations of $V$.
No intersection property or independence is assumed.
\end{lemma}
\begin{proof}
Some point in the $\gen{t_0,\ldots,t_{h-2}}$-orbit of $a$ is
moved by $t_{h-1}$. Otherwise that orbit would be fixed pointwise
by $t_{h-1}$ and invariant under $T$, hence equal to $V$, forcing
$t_{h-1}=1$. Choose a shortest word $w$ taking $a$ to a point
moved by $t_{h-1}$. Its point path is simple, and all earlier
vertices are fixed by $t_{h-1}$. The image of its endpoint under
$t_{h-1}$ is a further vertex. Thus $|w|\leq |V|-2\leq h$.

Every index $0,\ldots,h-2$ occurs in $w$. If index $q$ were
absent, the factors of smaller and larger index would commute.
The larger factors fix $a$, and the smaller factors commute
with $t_{h-1}$, which fixes $a$. The endpoint would consequently
be fixed as well. The first occurrences of the indices are in
increasing order: at the first failure, a new index at least
two greater than all previous indices commutes past the whole
preceding word and fixes $a$, so can be deleted, contradicting
minimality. The word therefore has length $h-1$ or $h$.

In the latter case exactly one index $q$ occurs twice. Between
its two occurrences there must be a new index; otherwise they
are consecutive and cancel. These intervening new indices,
in increasing order, end at some index $p\geq q+1$. After the
second occurrence of $q$, all remaining letters have index at
least $p+1\geq q+2$. Also $q\leq h-3$, because there was a new
larger index between its occurrences. The second $t_q$ therefore
commutes with the remaining letters and with $t_{h-1}$. Deleting
it preserves the property that the endpoint is moved by
$t_{h-1}$, again contradicting minimality. Hence
$w=t_0t_1\cdots t_{h-2}$.

This gives the displayed path, including its last $t_{h-1}$-edge.
For $j<i$, the generator $t_i$ commutes with every letter taking
$a$ to $a_j$ and fixes $a$, so fixes $a_j$.

First suppose that these $h+1$ path vertices exhaust $V$.
Descending from the last index shows
\[
 \gen{t_i,\ldots,t_{h-1}}=\Sym(\{a_i,\ldots,a_h\}).     \tag{2}
\]
Indeed the last mark is a single transposition. At the next
step $t_i$ has its displayed transposition and can only act
further on $W=\{a_{i+2},\ldots,a_h\}$. It centralizes the
already identified full symmetric group on $W$. Its additional
restriction is trivial if $|W|\ne2$; when $|W|=2$ it may also
be that pair's transposition, which already belongs to the
later subgroup and can be removed. Thus the displayed adjacent
transposition is always in the generated group, proving (2)
and hence (1).

Now suppose there is one extra point $z$, and every $t_i$ is
odd. Descending again, every $t_i$ is exactly its displayed
path transposition. This is immediate for the last mark.
At the next step the only possible additional action exchanges
$a_h,z$, which would make $t_{h-2}$ even and is forbidden.
At any earlier step the later marks generate the full symmetric
group on $W=\{a_{i+2},\ldots,a_h\}$, of size at least two.
Centralization preserves this nontrivial orbit and its singleton
complement $\{z\}$ among the remaining points. The centralizer
on $W$ is trivial for $|W|\geq3$; for $|W|=2$ its nonidentity
element would again make $t_i$ even. Thus $z$ is fixed by every
mark, contrary to transitivity. This proves all claims.
\end{proof}

Retain now the full critical six-point-head hypotheses and all
paid predecessors for L044: $n=2r\geq18$, the actual head is
supported on $D$, and the critical tail is natural alternating.
In particular the entire proof of
\emph{Pure alternating kernels in the critical prefixes} supplies
the strictly growing head orbits $C_k$ and the actual formula
\[
 P_k=G_{\leq k}
 =\{(s,u)\in\Sym(C_k)\times J_k:
                          \operatorname{sgn}s=\operatorname{sgn}u\},
                                                        \tag{3}
\]
with $C_3=D\cup\{b\}$, $|C_3|=7$, $|C_4|=8$.

\begin{theorem}[The terminal first-shared case is impossible]\label{end}
If colour $r-1$ has exactly one eligible edge, its first shared
label cannot be $3$.
\end{theorem}
\begin{proof}
Assume the contrary. The two orbits of $G_{r-1}=P_{r-2}$ are
$C=C_{r-2}$ and its complement $A$: $C$ is its actual head orbit,
and the unique crossing must connect all omission orbits.
Write the unique crossing as $\{c,a\}$ with $c\in C$
and $a\in A$, and put $h=r-4\geq5$. The head fixes $A$, and
its transitive action is therefore generated by the restrictions
of $\rho_3,\ldots,\rho_{r-2}$. The first of these is nontrivial
on $A$, by the shared-label hypothesis. Every generator of index
at most $r-3$ fixes $a$: it commutes with $\rho_{r-1}$ and
preserves the two omission orbits, whereas the crossing is unique.
Apply Lemma~\ref{path} in the reversed order
\[
            t_i=\rho_{r-2-i}|_A\quad(0\leq i<h).
\]
The strict prefix growth gives $|C|\geq r+2=h+6$, so
$|A|\leq h+2$. The lemma also gives $|A|\geq h+1$.
There are precisely two cases.

Suppose first that $|A|=h+1$, so $|C|=h+7=r+3$.
Put $C'=C_{r-3}$ and $\Delta=C\setminus C'$. Strict growth
gives $1\leq|\Delta|\leq2$.
The group $P'=P_{r-3}$ fixes $a$ and preserves separately
$A\setminus\{a\}$ and $\Delta$. By (1), its image on
$A\setminus\{a\}$ is the full $S_h$.
Let $J'$ be its literal exterior image on
$\Delta\sqcup A$. Its commutator subgroup is trivial on
$\Delta$, since $S_2$ is abelian, and maps onto
$(S_h)'=A_h$ on $A\setminus\{a\}$.
It follows that $J'$ contains the actual pure alternating
group on $A\setminus\{a\}$, fixing $\Delta\cup\{a\}$.
These are even exterior permutations. Formula (3) therefore
puts that same pure $A_h$ in $P'$, fixing $C'$ as well.

The final generator $\rho_{r-1}$ centralizes $P'$ and hence
both of its natural alternating subgroups on $C'$ and
$A\setminus\{a\}$. Their degrees are at least four, so it
fixes both sets pointwise. It can move only
$\Delta\cup\{a\}$, of size at most three. An even involution
on at most three points is the identity, contradicting its
crossing edge.

Suppose instead that $|A|=h+2$. Then $|C|=h+6=r+2$ and every
strict increment from $C_3$ to $C_{r-2}$ has size one.
On $C$, each $\rho_k$ with $3\leq k\leq r-2$ is a single
transposition. Here are the details needed to exclude any
exterior decoration on this particular orbit.
For $k\geq5$, $\rho_k$ fixes $C_{k-2}$ pointwise. Strict
growth forces it to exchange the unique point of
$C_{k-1}\setminus C_{k-2}$ with the unique point of
$C_k\setminus C_{k-1}$; there are no other points in $C_k$
on which it could act. The same assertion holds for $k=3,4$
by their explicit seven- and eight-point core formulas.
When $k<r-2$, the unique point of $C_{k+1}\setminus C_k$
is fixed by $\rho_k$, since $\rho_k$ preserves $C_k$ and
$C_{k+1}$. Each subsequent new point is obtained from the
previous one by a generator of index at least $k+2$.
String commutation consequently makes $\rho_k$ fix all these
later points as well. This proves the assertion on all of $C$.

Each $\rho_k$ is even globally, so its restriction to $A$ is
odd. Thus every $t_i$ in the displayed anchored action on
$h+2$ points is odd, contradicting the final assertion of
Lemma~\ref{path}. The two cases are exhausted.
\end{proof}

\paragraph{Attribution and boundary.}
The critical-head, natural-tail and prefix-kernel arguments remain
complete acquired dependencies, including all induction and finite
inputs. The last of these attributes the initial parity-fibre idea
to the root collaborator. The shortest-word argument with one
repeated label, the exact symmetric image on the exhausted path,
and the use of its derived subgroup to remove at most two exterior
points are the new proof here. The source route is the exceptional
first-head exit in L044 Propositions 5.26 and 5.28, physical pages
34 and 36--37; the source endpoint strategy is retained in attribution.
Together with the separate interior-cut proof, this excludes first
shared index $3$ at every unique cut of index at least $5$.
Other first shared indices, opposite ends, the terminal singleton
and the full source contracts remain separate. No whole-original,
conjecture, numerical or formalization credit is claimed.
\end{document}
